\documentclass[12pt]{article}

\setlength{\headheight}{15pt}
\usepackage{graphicx}		
\usepackage{amsmath}			
\usepackage[margin=1in, headheight=7pt, headsep=12pt]{geometry}
\usepackage{fancyhdr}		
\usepackage{enumerate}          
\usepackage[shortlabels]{enumitem}
\usepackage{amssymb}
\usepackage{amsmath}



\pagestyle{fancy}
\fancyhead[l]{\textbf{MATH 101}}
\fancyhead[c]{\textbf{AY 2026-2027 1st Sem}}
\fancyhead[r]{\textbf{Midterms Mock Exam}}
\fancyfoot[c]{\thepage}
\renewcommand{\headrulewidth}{0.1pt} 
\setlength{\headheight}{12pt}

\begin{document}

\noindent \textbf{PART I: MULTIPLE CHOICE QUESTIONS (20 Items — 2 Points Each)}
\vspace{0.5em}

\noindent \textit{Select the best answer for each question.}
    \begin{enumerate}[1.]
        \item{Which of the following is a valid proposition?}
            \begin{enumerate}[A.]
                \item{"Is the set of all real numbers uncountable?"}
                \item{"Solve the equation $x^2 - 5x + 6 = 0$."}
                \item{"For every prime number $p$, $p+1$ is composite."}
                \item{"$x + y = z$ for some $x, y \in \mathbb{R}$."}
            \end{enumerate}
            
        \item{Consider the argument: "If a function is differentiable, then it is continuous. Function $f$ is continuous. Therefore, $f$ is differentiable." What logical error or fallacy is demonstrated here?}
            \begin{enumerate}[A.]
                \item{Ad Hominem}
                \item{Appeal to Ignorance}
                \item{Fallacy of Affirming the Consequent}
                \item{Fallacy of Denying the Antecedent}
        \end{enumerate}
        
        \item{What is the negation of the conditional statement $p \to (q \vee \neg r)$?}
            \begin{enumerate}[A.]
                \item{$\neg p \vee (\neg q \wedge r)$}
                \item{$p \wedge \neg q \wedge r$}
                \item{$\neg p \to (\neg q \wedge r)$}
                \item{$p \to (\neg q \wedge r)$}
            \end{enumerate}
            
        \item{Let $p \to q$ be a conditional statement. Which of the following is logically equivalent to its inverse?}
        \begin{enumerate}[A.]
                \item{The original statement $p \to q$}
                \item{The contrapositive $\neg q \to \neg p$}
                \item{The converse $q \to p$}
                \item{$\neg p \wedge q$}
            \end{enumerate}
            
        \item{Let $P(x, y)$ be the predicate "$x + y = 0$" over the domain of real numbers $\mathbb{R}$. Which of the following statements is FALSE?}
        \begin{enumerate}[A.]
                \item{$\forall x \, \exists y \, P(x, y)$}
                \item{$\exists y \, \forall x \, P(x, y)$}
                \item{$\exists x \, \exists y \, P(x, y)$}
                \item{$\forall x \, \forall y \, \neg P(x, y)$ is false}
            \end{enumerate}
\newpage            
        \item{What is the negation of the quantified statement $\forall x \in \mathbb{R}, \, (x > 0 \to \exists y \in \mathbb{R} \, (xy = 1))$?}
        \begin{enumerate}[A.]
                \item{$\exists x \in \mathbb{R}, \, (x \le 0 \wedge \forall y \in \mathbb{R} \, (xy \neq 1))$}
                \item{$\exists x \in \mathbb{R}, \, (x > 0 \wedge \forall y \in \mathbb{R} \, (xy \neq 1))$}
                \item{$\forall x \in \mathbb{R}, \, (x > 0 \wedge \exists y \in \mathbb{R} \, (xy \neq 1))$}
                \item{$\exists x \in \mathbb{R}, \, (x > 0 \to \forall y \in \mathbb{R} \, (xy \neq 1))$}
            \end{enumerate}
            
        \item{Suppose you want to prove the statement: "If $3n + 2$ is odd, then $n$ is odd" using Proof by Contraposition. What should your starting assumption be?}
        \begin{enumerate}[A.]
                \item{Assume $3n + 2$ is even.}
                \item{Assume $n$ is odd.}
                \item{Assume $n$ is even.}
                \item{Assume $3n + 2$ is odd and $n$ is even.}
            \end{enumerate}
            
        \item{To disprove the statement $\forall x \in \mathbb{R}, \, x^2 > x$, which of the following serves as a valid counterexample?}
            \begin{enumerate}[A.]
                \item{$x = -2$}
                \item{$x = 2$}
                \item{$x = \frac{1}{2}$}
                \item{Both A and C}
            \end{enumerate}
            
        \item{In a proof by Principle of Mathematical Induction (PMI) showing that $n! > 2^n$ for $n \ge k_0$, what is the correct base case $n_0$?}
            \begin{enumerate}[A.]
                \item{$n_0 = 1$}
                \item{$n_0 = 3$}
                \item{$n_0 = 4$}
                \item{$n_0 = 0$}
            \end{enumerate}
            
        \item{Why does Russell's Paradox ($R = \{X \mid X \notin X\}$) arise in Naive Set Theory?}
            \begin{enumerate}[A.]
                \item{Because sets are allowed to be infinite.}
                \item{Because unrestricted set comprehension allows self-referential definitions without a bounding set.}
                \item{Because the Axiom of Choice fails for uncountable sets.}
                \item{Because the empty set is a subset of every set.}
            \end{enumerate}
            
        \item{Which Zermelo-Fraenkel (ZF) axiom guarantees that for any set $A$, the collection of all subsets of $A$ forms a set?}
            \begin{enumerate}[A.]
                \item{Axiom Schema of Separation}
                \item{Axiom of Pairing}
                \item{Axiom of Power Set}
                \item{Axiom Schema of Replacement}
            \end{enumerate}
            
        \item{Let $A = \{ \emptyset, \{\emptyset\} \}$. What is $|\mathcal{P}(A)|$?}
            \begin{enumerate}[A.]
                \item{2}
                \item{4}
                \item{8}
                \item{16}
            \end{enumerate}
            
        \item{Let $A = \{1, 2, 3\}$ and $B = \{2, 3, 4\}$. What is the symmetric difference $A \Delta B = (A - B) \cup (B - A)$?}
            \begin{enumerate}[A.]
                \item{$\{2, 3\}$}
                \item{$\{1, 4\}$}
                \item{$\{1, 2, 3, 4\}$}
                \item{$\emptyset$}
            \end{enumerate}
            
        \item{If $|A| = 3$ and $|B| = 4$, how many elements are in the set $(A \times B) \cap (B \times A)$?}
            \begin{enumerate}[A.]
                \item{Always 12}
                \item{Always 0}
                \item{$|A \cap B|^2$}
                \item{$|A \cup B|^2$}
            \end{enumerate}
            
        \item{Let $A = \{1\}$, $B = \{2\}$, and $C = \{3\}$. Which of the following is an element of the Cartesian product $(A \times B) \times C$?}
            \begin{enumerate}[A.]
                \item{$(1, 2, 3)$}
                \item{$((1, 2), 3)$}
                \item{$(1, (2, 3))$}
                \item{$\{\{1, 2\}, 3\}$}
            \end{enumerate}
            
        \item{Let $A, B, C,$ and $D$ be sets. Suppose that $A \times B \subseteq C \times D$. Which of the following conditions guarantees that $A \subseteq C$?}
            \begin{enumerate}[A.]
                \item{$B \neq \emptyset$}
                \item{$A \neq \emptyset$}
                \item{$C \neq \emptyset$}
                \item{$B \subseteq D$}
            \end{enumerate}
            
        \item{Under the definition of an ordered pair, $(a,b)=\{\{a\},\{a,b\}\}$. What is the simplified set representation of $(x,x)$?}
            \begin{enumerate}[A.]
                \item{$\{\{x\}\}$}
                \item{\{\{x\},\{x\}\}}
                \item{$\{\emptyset, \{x\}\}$}
                \item{$\{\{x\}, \emptyset\}$}
        \end{enumerate}
        
        \item{Let $A = \{x \in \mathbb{R} \mid x^2 - 3x + 2 = 0\}$ and $B = \{y \in \mathbb{Z} \mid 0 \le y < 2\}$. What is the cardinality of the power set of $A \times B$, denoted $\vert{}\mathcal{P}(A \times B)\vert{}$?}
            \begin{enumerate}[A.]
                \item{4}
                \item{8}
                \item{16}
                \item{32}
            \end{enumerate}

        \item{Which of the following set identities holds true for all sets $A, B,$ and $C$?}
            \begin{enumerate}[A.]
                \item{$A \times B = B \times A$}
                \item{$(A \times B) \times C = A \times (B \times C)$}
                \item{$A \times (B - C) = (A \times B) - (A \times C)$}
                \item{$(A \times B) \cup (B \times A) = (A \cup B) \times (A \cup B)$}
        \end{enumerate}
        \item{For any sets $A, B, C,$ and $D$, the expression $(A \times B) \cap (C \times D)$ is always equal to:}
            \begin{enumerate}[A.]
                \item{$(A \cap C) \times (B \cap D)$}
                \item{$(A \cap B) \times (C \cap D)$}
                \item{$(A \cup C) \times (B \cup D)$}
                \item{$(A \times C) \cap (B \times D)$}
            \end{enumerate}
    \end{enumerate}

\begin{center}
  \text{--- END OF PART I ---}
\end{center}
\newpage

\noindent \textbf{PART II: PROBLEM-SOLVING AND FORMAL PROOFS (11 Items — 5 Points Each)}
\begin{enumerate}[1.]
    \item{\textbf{Truth Tables \& Logical Equivalences}:}
    
    Without using truth tables, construct a step-by-step proof using standard logical equivalences (De Morgan's, Distributive, Implication laws, etc.) to show that:
    \[((p \to q) \wedge (q \to r)) \to (p \to r) \equiv \text{Tautology}\]

    \item{\textbf{Method of Deduction / Argument Validity:}}

    Determine whether the following argument is valid. If valid, provide a formal proof of validity using rules of inference (Modus Ponens, Modus Tollens, Hypothetical Syllogism, Disjunctive Syllogism, etc.). If invalid, construct a truth assignment counterexample.
    \begin{itemize}
        \item{Premise 1: $\neg p \vee q$}
        \item{Premise 2: $r \to \neg q$}
        \item{Premise 3: $r \vee s$}
        \item{Premise 4: $\neg s$}
        \item{Conclusion: $\neg p$}
    \end{itemize}

    \item{\textbf{Proof by Contraposition}:}

    Let $x \in \mathbb{R}$. Prove by contraposition: If $x^3 + 2x^2 + x < 0$, then $x < 0$.

    \item{\textbf{Proof by Contradiction}:}

    Prove that if $n \in \mathbb{Z}^+$, then $n^2 + 2$ is not divisible by 4.

    \item{\textbf{Proof by Cases}:}

    Prove that for all integers $n$, $n^3 - n$ is always divisible by 6.

    \item{\textbf{Principle of Mathematical Induction (PMI)}:}
        \begin{enumerate}[a.]
            \item{Use mathematical induction to prove that $5^{2n} - 1$ is divisible by 24 for all $n \ge 1$.} 
            \item{Use PMI to prove that for all integers $n \ge 1$:}

            $$1 \cdot 2 + 2 \cdot 3 + 3 \cdot 4 + \dots + n(n+1) = \frac{n(n+1)(n+2)}{3}$$
        \end{enumerate}
    

    \item{\textbf{Set Operations \& Formal Set Proof:}}

    Let $A, B,$ and $C$ be sets inside a universal set $U$. Prove element-wise that:
    \[A - (B \cup C) = (A - B) \cap (A - C)\]

    \item{\textbf{Cartesian Product:}}
    \begin{enumerate}[a.]
    \vspace{-0.5em}
        \item{Let $A$ and $B$ be non-empty sets such that $A \times B = B \times A$. Prove that $A = B$.}
        \item{Let $A, B, C,$ and $D$ be sets. Prove that $(A \times B) \cap (C \times D) = (A \cap C) \times (B \cap D)$.}
        \item{Prove that for any sets $A, B,$ and $C$, $A \times (B - C) = (A \times B) - (A \times C)$.}
    \end{enumerate}
    
\end{enumerate}
\begin{center}
  \text{--- END OF PART II ---}
\end{center}

\text{References: https://bit.ly/MATH101Handouts}
\end{document}